\documentclass[11pt,a4paper]{article}

\usepackage[utf8]{inputenc}
\usepackage[T1]{fontenc}
\usepackage[margin=2cm]{geometry}
\usepackage{graphicx, booktabs, amsmath, hyperref, float, adjustbox, underscore}
\hypersetup{colorlinks=true, linkcolor=blue, urlcolor=blue}

\title{Hantavirus Data -- Descriptive Statistics Analysis}
\author{}
\date{}

\begin{document}
\maketitle
\thispagestyle{empty}

\section{Introduction}

This project applies descriptive statistics concepts learned from a textbook to a real
hantavirus dataset from Kaggle, as a practice project.

Data source: \texttt{https://www.kaggle.com/datasets/zkskhurram/hantavirus-andes-virus-global-epidemiology}

\noindent\textbf{Disclaimer:} This is a historical Kaggle dataset on hantavirus epidemiology, not the 2025 outbreak data. This project focuses on learning descriptive statistics, not on current events.

\section{Data Cleaning}

No data errors were found. Clinical symptoms have structural missingness by syndrome.
All outliers (Table~\ref{tab:outliers}) were genuine extreme values, not errors.

\section{Univariate Distribution Analysis}

\subsection{Incubation Days}
Right-skewed (skew=0.47), unimodal, thin tails. Mean=20.4d, Median=19d.
Not normal. Incubation differs by syndrome: HPS (median=16d) vs HFRS (median=22d).

\subsection{Hospital Days}
Appears bimodal when aggregated -- an artifact of mixing four severity groups with
distinct non-overlapping centers (Mild=2d, Moderate=6d, Severe=13d, Critical=21d).

\subsection{ICU Days}
Spike-and-slab: 50\% never admitted. Mild/Moderate never go to ICU;
Severe/Critical always do. Among admitted: median=5d.

\subsection{Environmental Variables}
Temperature is approximately normal (skew=-0.15). Rainfall is right-skewed (skew=1.80),
10x variation by biome. Rodent abundance is symmetric and constant across biomes.

\subsection{Confirmed Cases}
Extreme right skew (skew=2.82). Mean=1793, Median=58. Tail driven by China HFRS.

\subsection{Case Fatality Rate}
Bimodal by syndrome. HPS (28.1\%) is 20x more lethal than HFRS (1.4\%).

\section{Comparative Analysis}
Incubation is similar across severities but differs by syndrome. Hospital/ICU days
separate cleanly by severity. Cases and CFR differ dramatically by syndrome.

\section{Key Takeaways}
\begin{enumerate}
\item Check subgroups before describing a distribution
\item For skewed data: median + IQR, not mean + std
\item Spike-and-slab distributions need split analysis
\item Mixing groups creates misleading aggregates
\item Bar charts with error bars hide the truth
\end{enumerate}

\section{Tables}

\begin{table}[htbp]
\centering
\resizebox{\textwidth}{!}{
\begin{table}[htbp]
\caption{Summary statistics for all numeric variables across datasets}
\label{tab:master}
\begin{tabular}{lrlllllllll}
\toprule
Variable & n & Mean & Median & Std & IQR & MAD & Skewness & Kurtosis & Best Center & Best Spread \\
\midrule
Incubation Days & 7538 & 20.41 & 19.00 & 8.07 & 12.00 & 5.00 & 0.47 & -0.49 & Mean & Std \\
Hospital Days & 7538 & 10.46 & 9.00 & 6.18 & 9.00 & 4.00 & 0.60 & -0.42 & Median & Iqr \\
ICU Days & 7538 & 3.30 & 0.00 & 4.16 & 5.00 & 0.00 & 1.19 & 0.50 & Median & Iqr \\
Temperature (°C) & 896 & 11.26 & 11.70 & 7.18 & 9.70 & 4.90 & -0.15 & -0.25 & Mean & Std \\
Rainfall (mm) & 896 & 252.83 & 195.35 & 217.76 & 237.10 & 108.60 & 1.80 & 5.05 & Median & Iqr \\
Rodent Abundance & 896 & 0.67 & 0.67 & 0.18 & 0.30 & 0.15 & 0.02 & -1.12 & Mean & Std \\
Confirmed Cases & 939 & 1792.77 & 58.00 & 4470.37 & 415.00 & 51.00 & 2.82 & 7.14 & Median & Iqr \\
Case Fatality Rate & 939 & 0.09 & 0.02 & 0.13 & 0.17 & 0.02 & 1.28 & 0.08 & Median & Iqr \\
\bottomrule
\end{tabular}
\end{table}

}
\end{table}

\begin{table}[htbp]
\centering
\resizebox{\textwidth}{!}{
\begin{table}[htbp]
\caption{Hospital days stratified by severity level — each group has a distinct, non-overlapping center}
\label{tab:hospital_severity}
\begin{tabular}{lrllll}
\toprule
Severity & n & Mean & Median & Std & IQR \\
\midrule
Mild & 1121 & 2.5 & 2 & 0.5 & 1.0 \\
Moderate & 2669 & 6.5 & 6 & 1.2 & 3.0 \\
Severe & 2606 & 13.5 & 13 & 2.4 & 5.0 \\
Critical & 1142 & 20.6 & 21 & 3.6 & 7.0 \\
\bottomrule
\end{tabular}
\end{table}

}
\end{table}

\begin{table}[htbp]
\centering
\resizebox{\textwidth}{!}{
\begin{table}[htbp]
\caption{ICU admission rates and stay lengths by severity}
\label{tab:icu_severity}
\begin{tabular}{lrlll}
\toprule
Severity & n & ICU Rate & ICU Days (median) & ICU Days (mean) \\
\midrule
Mild & 1121 & 0\% & — & — \\
Moderate & 2669 & 0\% & — & — \\
Severe & 2606 & 100\% & 4 & 4.5 \\
Critical & 1142 & 100\% & 12 & 11.5 \\
\bottomrule
\end{tabular}
\end{table}

}
\end{table}

\begin{table}[htbp]
\centering
\resizebox{\textwidth}{!}{
\begin{table}[htbp]
\caption{Case fatality rate by syndrome — HPS is ~20x more lethal than HFRS}
\label{tab:cfr_syndrome}
\begin{tabular}{lrllll}
\toprule
Syndrome & n & Mean CFR & Median CFR & Max CFR & IQR \\
\midrule
HPS & 273 & 28.1\% & 29.0\% & 46.2\% & 12.3\% \\
HFRS & 666 & 1.4\% & 0.6\% & 6.5\% & 2.9\% \\
\bottomrule
\end{tabular}
\end{table}

}
\end{table}

\begin{table}[htbp]
\centering
\resizebox{\textwidth}{!}{
\begin{table}[htbp]
\caption{Confirmed cases by WHO region — WPRO (Asia) dominates}
\label{tab:cases_region}
\begin{tabular}{lrllll}
\toprule
WHO Region & n & Mean & Median & Max & Total \\
\midrule
AMRO & 273 & 30 & 13 & 146 & 8111 \\
WPRO & 227 & 5358 & 482 & 23885 & 1216301 \\
EURO & 439 & 1046 & 89 & 9496 & 459003 \\
\bottomrule
\end{tabular}
\end{table}

}
\end{table}

\begin{table}[htbp]
\centering
\resizebox{\textwidth}{!}{
\begin{table}[htbp]
\caption{Outlier investigation results — all flagged outliers were genuine extreme values, not errors}
\label{tab:outliers}
\begin{tabular}{lrll}
\toprule
Variable & n outliers & Percent & Verdict \\
\midrule
incubation_days & 0 & 0.0\% & — \\
hospital_days & 0 & 0.0\% & — \\
icu_days & 432 & 5.7\% & Keep (real) \\
avg\_temp\_c & 1 & 0.1\% & Keep (real) \\
rainfall\_mm & 44 & 4.9\% & Keep (real) \\
rodent\_abundance & 0 & 0.0\% & — \\
confirmed\_cases & 170 & 18.1\% & Keep (real) \\
case\_fatality\_rate & 9 & 1.0\% & Keep (real) \\
\bottomrule
\end{tabular}
\end{table}

}
\end{table}

\end{document}
